% from ega4-1-fr.tex
\oldpage[IV-1]{241}

\begin{proposition}[1.8.7]
\label{IV.1.8.7-fr}
Soient $S$ un préschéma, $g:X\to S$, $h:Y\to S$ deux morphismes de
présentation finie, $f:X\to Y$ un $S$-morphisme. Pour tout $s\in S$, posons
$X_s=g^{-1}(s)=X\times_S\operatorname{Spec}(k(s))$,
$Y_s=h^{-1}(s)=Y\times_S\operatorname{Spec}(k(s))$,
$f_s=f\times 1:X_s\to Y_s$. Alors l'ensemble des $s\in S$ tels que $f_s$
soit surjectif (resp. radiciel) est localement constructible.
\end{proposition}

\begin{proof}
Soit $E$ l'ensemble des $s\in S$ tels que $f_s$ soit surjectif~; l'ensemble
$S-E$ est égal à $h(Y-f(X))$~; or $f$ est de présentation finie (1.6.2, (v))
donc $f(X)$ est localement constructible dans $Y$ (1.8.4)~; comme $h$ est de
présentation finie, $h(Y-f(X))$ est localement constructible dans $S$, donc
aussi $E$.

Pour montrer que l'ensemble $E'$ des $s\in S$ tels que $f_s$ soit radiciel
est localement constructible, nous utiliserons le
\end{proof}

\begin{lemma}[1.8.7.1]
\EGAReaderTarget{ega4.c04.lem.1-8-7-1}{IV.1.8.7.1-fr}
Soient $U$, $V$ deux préschémas~; pour qu'un morphisme $f:U\to V$ soit
radiciel, il faut et il suffit que le morphisme diagonal
$\Delta_f:U\to U\times_VU$ soit surjectif. Par suite, tout morphisme radiciel
est séparé.
\end{lemma}

\begin{proof}
En effet, $\Delta_f$ étant une immersion (I, 5.3.9), est un morphisme
localement de type fini (1.3.4)~; dire que $\Delta_f$ est surjectif signifie
donc que pour tout corps algébriquement clos $K$, l'application
correspondante $U(K)\to(U\times_VU)(K)=U(K)\times_{V(K)}U(K)$ est surjective
(1.3.7 et I, 3.4.2.1)~; mais par définition d'un produit fibré d'ensembles,
cela signifie que l'application $U(K)\to V(K)$ correspondant à $f$ est
injective, et cela équivaut à dire que $f$ est radiciel (I, 3.5.5).

Cela étant, $\Delta_{f_s}$ se déduit de $\Delta_f:X\to X\times_YX$ par le
changement de base $\operatorname{Spec}(k(s))\to S$ (I, 5.3.4). Comme $f$
est de présentation finie, il en est de même du morphisme structural
$X\times_YX\to Y$ (1.6.2, (iv)), donc aussi de $\Delta_f$ (1.6.2, (v))~; il
suffit donc d'appliquer la première partie de la proposition à $\Delta_f$,
en utilisant le lemme (1.8.7.1).
\end{proof}

\subsection{Ensembles pro-constructibles et ind-constructibles}
\label{subsection:IV.1.9-fr}

\begin{lemma}[1.9.1]
\label{IV.1.9.1-fr}
Soit $(A_\alpha)_{\alpha\in I}$ un système inductif d'anneaux dont l'ensemble
d'indices $I$ est filtrant à droite, et soit $A=\varinjlim A_\alpha$. Pour que
$A=0$, il faut et il suffit qu'il existe un indice $\gamma$ tel que
$A_\gamma=0$ (et l'on a alors $A_\beta=0$ pour tout $\beta\geq\gamma$).
\end{lemma}

\begin{proof}
En effet, pour tout $\alpha\in I$, l'homomorphisme canonique
$\varphi_\alpha:A_\alpha\to A$ transforme l'élément unité en élément unité~;
dire que $A=0$ signifie que $\varphi_\alpha(1)=0$, ou encore
$\varphi_\alpha(1)=\varphi_\alpha(0)$~; on sait que cela équivaut à dire
qu'il existe un indice $\gamma\geq\alpha$ tel que l'homomorphisme
$\varphi_{\gamma\alpha}:A_\alpha\to A_\gamma$ soit tel que
$\varphi_{\gamma\alpha}(1)=\varphi_{\gamma\alpha}(0)$, et cette dernière
relation équivaut à $A_\gamma=0$, d'où le lemme.
\end{proof}

\begin{proposition}[1.9.2]
\label{IV.1.9.2-fr}
Soient $B$ un anneau, $(A_\alpha)_{\alpha\in I}$ un système inductif de
$B$-algèbres dont l'ensemble d'indices $I$ est filtrant à droite, et soit
$A=\varinjlim A_\alpha$~; posons $Y=\operatorname{Spec}(B)$,
$X_\alpha=\operatorname{Spec}(A_\alpha)$, $X=\operatorname{Spec}(A)$, et
soient $u:X\to Y$, $u_\alpha:X_\alpha\to Y$ les morphismes correspondant aux
homomorphismes structuraux $B\to A$, $B\to A_\alpha$. Alors~:
\end{proposition}

\begin{env}[1.9.2.i]
\label{IV.1.9.2.i-fr}
\emph{Pour que $X=\varnothing$, il faut et il suffit qu'il existe un
$\gamma\in I$ tel que $X_\gamma=\varnothing$, et l'on a alors
$X_\alpha=\varnothing$ pour $\alpha\geq\gamma$.}
\end{env}

